% LoG 2026 version. Build modes:
%   \usepackage[review]{log_2026}    -> anonymous (OpenReview submission)
%   \usepackage[preprint]{log_2026}  -> de-anonymized preprint
%   \usepackage{log_2026}            -> camera-ready
\documentclass{article}
\usepackage[review]{log_2026}
\usepackage{booktabs}
\usepackage{multirow}
\usepackage{amsfonts}
\usepackage{amsmath}
\usepackage{graphicx}
\usepackage[numbers,compress,sort]{natbib}

\title[How Much of the Elliptic Leaderboard Is Real?]{How Much of the Elliptic
Leaderboard Is Real? Decomposing Evaluation Inflation into Leakage, Shift, and
Base-Rate Effects on Temporal Transaction Graphs}

\author[Audibert]{%
Julien Audibert\\
DeepFi, France\\
\email{ja@deepfi.io}
}

\begin{document}

\maketitle

\begin{abstract}
The Elliptic Bitcoin dataset and its extensions are standard benchmarks for
graph-based fraud detection, and recent work reports illicit-detection scores
approaching 0.98. We show that these numbers are largely an artifact of the
evaluation protocol. Holding models, features, and training fixed, a random split
inflates illicit-class PR-AUC by $+0.195$ (XGBoost) to $+0.553$ (GCN) over five
paired seeds, relative to the standard temporal split. Moreover, the full
published recipe, a random split combined with aggregate metrics, reproduces the
reported headline numbers with a generic XGBoost (accuracy 0.992, weighted F1
0.992). We then decompose this gap with two controls. First, a prevalence-matched
random test set isolates a small base-rate component. Second, an inductive
ablation, which removes all test nodes from the graph during training, shows that
the message-passing leakage we initially hypothesized is indistinguishable from
zero for all four GNN architectures tested, including a graph transformer. The dominant component is
\emph{distribution access}: a random split trains the model on the test period's
distribution, and the models that generalize worst across temporal shift gain the
most. Consistent with this mechanism, the inflation tracks the drift
magnitude across four datasets (stable $+0.001$, moderate drift $+0.124$, regime
change $+0.41$ to $+0.55$) and concentrates in the post-shift window. Under the
temporal protocol, gradient-boosted trees still lead graph neural networks, and
four drift-robust variants, including a first adaptation of USAD to graphs, fail
to improve on the baselines. Every reported number can be regenerated with a
single seeded command from our repository.
\end{abstract}

\section{Introduction}
Detecting illicit activity in blockchain transaction graphs is a canonical
application of graph machine learning. The Elliptic dataset
\citep{weber2019aml} (203{,}769 Bitcoin transactions across 49 time steps,
$\sim$2\% illicit) and Elliptic++ \citep{elmougy2023elliptic} (adding 822{,}942
wallet addresses) are the de facto benchmarks. A stream of recent papers reports
illicit-detection scores approaching 0.98, for example accuracy 0.9802 and
weighted F1 0.9799 under a stratified random split \citep{lin2026wavelet}, or
random-split evaluations on Elliptic++ \citep{nguyen2025sagefin}. This suggests
the problem is nearly solved. We argue this impression is false and traces to two
evaluation choices: (i) random rather than temporal train/test splits, which give
models access to the test period's distribution, and (ii) reporting weighted or
overall F1, dominated by the $\sim$98\% licit class, rather than illicit-class
metrics.

This question matters to the graph learning community specifically. Elliptic is
one of the most cited node-classification benchmarks with real-world stakes, and
the gap we quantify is largest precisely for GNNs: the models this community
builds gain the most from the broken protocol. An evaluation regime that
systematically favors a model family distorts the apparent state of the art and
the direction of methods research.

Our contributions are: (1) a systematic, multi-model, multi-seed quantification of
the random-split inflation on Elliptic-family benchmarks; (2) a decomposition of
that inflation into base-rate, graph-leakage, and distribution-access components,
including the measured refutation of our own message-passing-leakage hypothesis;
(3) a temporally stable cross-domain control consistent with the shift-driven
mechanism; and (4) negative results, reported in full, for four drift-robust
variants. These negatives indicate that the post-shift regime is the genuinely
open problem, rather than an unexplored one. We release an evaluation harness in
which every reported number is one seeded command.

\section{Related work}
\paragraph{Evaluation integrity.} Leakage was formalized by \citet{kaufman2011leakage}
(the ``no time machine'' rule); \citet{kapoor2023leakage} taxonomize it as a driver
of ML reproducibility failures; \citet{arp2022dos} name ``temporal snooping'' among
security-ML pitfalls; \citet{shchur2018pitfalls} document protocol sensitivity in
GNN evaluation generally. The closest template is TESSERACT
\citep{pendlebury2019tesseract}, which showed Android-malware F1 of $\sim$0.99
collapsing under time-aware evaluation. Concurrently, \citet{roth2026leakage}
quantifies four leakage classes across 2{,}047 \emph{iid} tabular benchmarks,
finding normalization leakage negligible and boundary effects invisible under
random cross-validation. This is precisely the regime where our mechanism
predicts no inflation. Our results show the complementary regime: on temporally
drifting fraud benchmarks the split protocol becomes the dominant effect, through
distribution access rather than any classical leakage class. Relative to this
line, our contribution is graph-native and mechanistic: we quantify the inflation
on Elliptic-family benchmarks across model classes and decompose it, showing the
graph-specific channel is null and the dominant term is distribution access.

\paragraph{Temporal graph benchmarks.} The Temporal Graph Benchmark
\citep{huang2023tgb} standardizes chronological splits for temporal graph
learning tasks, and Chartalist \citep{shamsi2022chartalist} provides labeled
blockchain graph datasets with timestamps that make time-aware splits possible.
Our results supply a quantified argument for why this discipline is not optional
on fraud graphs: we measure what is gained, artificially, when it is dropped, and
show that the gain is largest for the models whose temporal generalization is
weakest.

\paragraph{Elliptic-specific evaluation.} \citet{weber2019aml} established the
temporal protocol and the trees-over-GNNs result; GADBench \citep{tang2023gadbench}
confirms tuned tree ensembles lead on Elliptic (tuned RF AUPRC $\approx$0.789,
matching our XGBoost). Closest to our ablation, a recent preprint
\citep{maganti2026liability} reports a large ($\sim$39.5 F1 point)
transductive-versus-inductive protocol gap on Elliptic under temporal evaluation:
training-time exposure to test-period graph structure inflates results. This is
consistent with our decomposition, which offers a candidate explanation rather
than standing in tension with our null: under a \emph{temporal} protocol,
test-period nodes are the model's only access to the shifted test distribution,
so graph visibility carries exactly the distributional information that drives
inflation; under our \emph{random}-split ablation, the labeled training set
already spans all periods, test-node features are redundant, and removing them
changes nothing ($-0.000\pm0.008$). On this reading the graph channel is not a
separate leak: it matters precisely insofar as it transports
otherwise-unavailable distributional information. In other words, a single
mechanism expresses itself under two protocols. A direct test (our inductive
ablation run under the temporal protocol) would adjudicate this reading and is
left to future work.

\paragraph{Concurrent work: construction-level leakage.} While this study was
underway, \citet{safar2026opacity} published a forensic deconstruction of the
Elliptic dataset itself: reverse-engineering over 90\% of its feature semantics,
they show that the ``local'' features embed address-level global statistics
computed at a reference time \emph{after} the test period (look-ahead by
construction), that feature normalization was fit on the full dataset, and that
shared addresses couple training and test samples through both features and
inherited labels. Their axis is complementary to ours: they identify and formalize
\emph{construction-level} leakage without quantifying its performance impact; we
quantify and decompose \emph{protocol-level} inflation without touching feature
construction. The two compose: both arms share the same features, so
\emph{exposure} to construction-level leakage is symmetric, though its induced
\emph{effect} need not be identical across arms; we therefore treat it as an
unmeasured layer that makes our temporal-arm scores \emph{upper bounds} on
deployable performance and could perturb individual components (flagged in
Limitations), while leaving the qualitative decomposition intact. The true
deployment-realistic numbers are therefore lower still, which strengthens rather
than weakens the overall conclusion. None of the works above measure random-split
inflation across models, decompose it into components, or test drift-graded
controls.

\paragraph{Graph anomaly detection and USAD.} Reconstruction autoencoders
\citep{ding2019dominant}, generative and adversarially regularized variants, and
one-class methods form the unsupervised graph-anomaly literature. Among
adversarial reconstruction schemes, we selected USAD \citep{audibert2020usad}, a
two-decoder autoencoder for multivariate time series, as the representative
approach for our negative-result study: it is explicitly designed for training
stability over GAN-style alternatives, which makes it a strong candidate under a
regime where training instability could otherwise explain a failure. To our
knowledge it had not previously been adapted to graphs; we adapt it and report
the outcome.

\section{Experimental protocol}
\label{sec:protocol}

\begin{table}[t]
\centering
\caption{Datasets. Drift regime is characterized in Section~\ref{sec:dose}:
Elliptic's labeled fraud prevalence collapses from 9.2\% to 2.5\% at the
dark-market shutdown (time step 43); AMLworld's laundering prevalence
doubles across the stream; DGraph-Fin's fraud rate is temporally stable.}
\label{tab:datasets}
\resizebox{\textwidth}{!}{%
\begin{tabular}{lllll}
\toprule
Dataset & Scale & Task & Split & Drift regime \\
\midrule
Elliptic \citep{weber2019aml} & 204k tx, 234k edges & tx classification & time step 34 & regime change \\
Elliptic++ \citep{elmougy2023elliptic} & +823k addresses, 2.87M edges & address classification & first-seen step & regime change \\
DGraph-Fin \citep{huang2022dgraph} & 3.7M nodes & node classification & first-seen time & stable \\
IBM AMLworld HI-Small & 5.08M tx (synthetic) & tx classification & transaction time & moderate \\
\bottomrule
\end{tabular}}
\end{table}

All experiments use one codebase, fixed seeds, and a CI-gated environment
(Table~\ref{tab:datasets}). The evaluation protocol is: a temporal split
(Elliptic's canonical early-versus-late time steps; first-seen splits for
addresses and DGraph-Fin), illicit-class PR-AUC and F1 as primary metrics, a
validation slice carved from the latest train steps, and decision thresholds
chosen on validation only. Leakage comparisons hold the model, features, and
trainer fixed and change only the split. Models: GCN \citep{kipf2017gcn},
GraphSAGE \citep{hamilton2017graphsage}, GAT \citep{velickovic2018gat}, a graph
transformer (two TransformerConv layers, the UniMP operator
\citep{shi2021unimp}, 4 heads), and XGBoost \citep{chen2016xgboost}. Unless stated otherwise, GNNs use two convolution
layers, hidden dimension 64, dropout 0.5, Adam with learning rate $10^{-2}$ and
weight decay $5\times10^{-4}$, up to 200 epochs with early stopping on
validation PR-AUC (patience 20), and a class-weighted cross-entropy loss; XGBoost
uses 300 trees of depth 6 with \texttt{scale\_pos\_weight} set from the train
class ratio. Features are standardized with a scaler fit on the training arm
only.

\paragraph{Statistics.} Multi-seed results use five paired seeds (42--46) unless
stated otherwise; we report mean $\pm$ sample standard deviation (ddof$=1$) and,
for headline claims, 95\% $t$-confidence intervals and exact sign-flip permutation
tests on the paired per-seed differences (the smallest attainable one-sided $p$ at
$n{=}5$ is $1/32$). The ``leakage is null'' statement is an equivalence claim and
is therefore assessed with two-one-sided tests (TOST): we report the smallest
round symmetric margin containing the 90\% CI; these margins are descriptive
(post hoc), not pre-registered. Per-seed values for every cell are released with
the code (Appendix~\ref{app:perseed}). The Elliptic++ address task, the
DGraph-Fin control, DANN, and the USAD variants use three paired seeds; the
per-architecture tuning sweeps, the EvolveGCN sweep, the DGraph width-128 probe,
and the step-2 full-train baseline are single runs at seed 42 and are labeled as
such where quoted.

\section{Results}
\subsection{Baseline ordering under the temporal protocol}
Under the temporal split, XGBoost (PR-AUC $0.791\pm0.002$; F1 0.82 at seed 42
under the full-train protocol) leads GraphSAGE ($0.510\pm0.039$), the graph
transformer ($0.487\pm0.038$), GAT ($0.297\pm0.061$) and GCN ($0.265\pm0.058$),
reproducing \citet{weber2019aml}. Transformer-style dot-product attention does
not close the temporal generalization gap: it lands within seed noise of
GraphSAGE, far below the tree baseline. A
12-configuration tuning sweep per architecture, selected on temporal validation
only (seed 42), does not change this picture: for GraphSAGE it improves F1
($0.422 \rightarrow 0.486$) while leaving PR-AUC within seed noise of the default
(0.506 vs $0.510\pm0.039$); for GCN, validation selection actually \emph{degrades}
test PR-AUC ($0.294 \rightarrow 0.274$), an instance of validation-test
misalignment under shift that we report rather than omit. Our GNN F1 remains below
some honest-protocol literature values ($\sim$0.69); we report this openly as a
protocol/training-detail discrepancy rather than claiming parity.

\subsection{The inflation, quantified}
\begin{table}[h]
\centering
\caption{Illicit-class PR-AUC under the temporal split versus a stratified
random split (mean $\pm$ std over 5 seeds). Same model, same trainer, same
features; only the split differs.}
\label{tab:inflation}
\begin{tabular}{lccc}
\toprule
Model & Temporal & Random & Inflation \\
\midrule
GCN & $0.265\pm0.058$ & $0.818\pm0.010$ & $+0.553\pm0.062$ \\
GraphSAGE & $0.510\pm0.039$ & $0.920\pm0.008$ & $+0.410\pm0.046$ \\
GAT & $0.297\pm0.061$ & $0.775\pm0.036$ & $+0.478\pm0.041$ \\
GraphTransformer & $0.487\pm0.038$ & $0.902\pm0.007$ & $+0.415\pm0.039$ \\
XGBoost & $0.791\pm0.002$ & $0.986\pm0.001$ & $+0.195\pm0.001$ \\
\bottomrule
\end{tabular}
\end{table}

\begin{figure}[t]
\centering
\includegraphics[width=0.75\linewidth]{figs/leakage_multi.png}
\caption{Temporal versus random-split PR-AUC per model (5 seeds). The
inflation is largest for the GNNs, the models whose temporal generalization is
weakest.}
\label{fig:leakage}
\end{figure}

Every inflation is significant at the smallest attainable exact permutation level
($p=1/32$), with 95\% CIs bounded away from zero (e.g.\ GraphSAGE $[0.353,0.467]$;
Figure~\ref{fig:leakage}). Under the random split, XGBoost reaches PR-AUC
$0.986\pm0.001$ (F1 0.955 at seed 42), inside the literature's reported
$\sim$0.90--0.98 band.

\paragraph{The published recipe, replicated exactly.} Following the full published
recipe, a stratified random split combined with aggregate metrics, a generic
XGBoost reproduces the headline numbers (5 seeds): accuracy
$0.9917\pm0.0002$ and weighted F1 $0.9916\pm0.0002$, matching and slightly
exceeding reported values such as accuracy 0.9802 / weighted F1 0.9799
\citep{lin2026wavelet}. We replicate the protocol class (split and metric
choices), not any specific pipeline's implementation; this shows that the
protocol alone suffices to reach the band. The same
model under the temporal protocol scores illicit-F1 $0.756\pm0.007$ / PR-AUC
$0.791\pm0.002$. Notably, \emph{accuracy remains 0.969 even under the temporal
split}: metric choice and split choice each contribute to the illusion, and
together they turn 0.76 into 0.99 with no novel architecture involved. The effect
generalizes to the Elliptic++ address task (heterogeneous graph, 822k address
nodes) under a deployment-honest feature protocol (pre-split wallet aggregation,
train-only scaling): over three seeds, PR-AUC $0.417\pm0.052$ (temporal) versus
$0.965\pm0.002$ (random), an inflation of $+0.548\pm0.050$.

\subsection{Decomposing the gap}
\label{sec:decompose}
Two controls, threshold-free PR-AUC, five seeds. (i) Matching the random test
set's prevalence to the temporal test's (6.5\%) isolates the base-rate component.
(ii) An inductive ablation removes all test nodes from the graph during training,
so no test feature or test-adjacent edge is visible; this measures graph-specific
message-passing leakage directly.

\begin{table}[h]
\centering
\caption{Decomposition of the temporal-versus-random gap (PR-AUC, 5 seeds).}
\label{tab:decompose}
\resizebox{\textwidth}{!}{%
\begin{tabular}{lccccc}
\toprule
Component & GCN & GraphSAGE & GAT & Transformer & XGBoost \\
\midrule
Base rate & $+0.061\pm0.004$ & $+0.026\pm0.006$ & $+0.073\pm0.013$ & $+0.033\pm0.002$ & $+0.005\pm0.001$ \\
Message-passing leakage & $+0.003\pm0.018$ & $-0.000\pm0.008$ & $-0.010\pm0.036$ & $+0.000\pm0.011$ & n/a \\
Distribution access & $+0.489\pm0.047$ & $+0.384\pm0.049$ & $+0.415\pm0.050$ & $+0.383\pm0.036$ & $+0.189\pm0.002$ \\
\midrule
Total gap & $+0.553\pm0.062$ & $+0.410\pm0.046$ & $+0.478\pm0.041$ & $+0.415\pm0.039$ & $+0.195\pm0.001$ \\
\bottomrule
\end{tabular}}
\end{table}

\begin{figure}[t]
\centering
\includegraphics[width=0.85\linewidth]{figs/decompose.pdf}
\caption{Decomposition of the temporal-versus-random PR-AUC gap. Bars are 5-seed
means; dots are individual seeds. The message-passing leakage component clusters
at zero for all four GNNs; distribution access dominates everywhere.}
\label{fig:decompose}
\end{figure}

The message-passing channel we ourselves hypothesized is \textbf{indistinguishable
from zero for all four architectures}: training with test nodes entirely removed
from the graph leaves random-split PR-AUC statistically unchanged. Because this is
an equivalence claim, we assess it with TOST: the 90\% CIs bound the effect within
$\pm 0.01$ for GraphSAGE ($[-0.008, +0.007]$), $\pm 0.02$ for the graph
transformer ($[-0.010, +0.010]$), $\pm 0.03$ for GCN ($[-0.014, +0.020]$), and
$\pm 0.05$ for GAT ($[-0.045, +0.025]$); the GAT bound is loose owing to its seed
variance, so we can rule out a large leak but not a small one there. No architecture shows a positive leak. Distribution access is
computed as the residual (total gap minus the base-rate and message-passing
components), so its uncertainty compounds theirs. The dominant term
everywhere is distribution access: a random split trains the model on examples
drawn from the test period, an advantage no deployed model can have. GNNs inflate
more than XGBoost (paired GraphSAGE$-$XGBoost inflation difference $+0.215$, 95\%
CI $[0.158, 0.272]$, $p=1/32$) because their temporal generalization is weaker,
not because of a graph-specific leak. We report this as an explicit
self-correction. This null is also consistent with the large
transductive-versus-inductive gap reported by \citet{maganti2026liability} under
temporal evaluation (Section 2), on the candidate reading that graph visibility
inflates results only when it is the sole conduit of test-period distributional
information; when distribution access is already granted by the split, the graph
channel adds nothing measurable.

\subsection{The mechanism's falsifiable prediction: inflation concentrates in the
shifted window}
\label{sec:windowed}
If distribution access drives the inflation, it should be largest where the
distribution shift is largest: the post-shutdown window (steps $\geq$43), where the
labeled fraud prevalence collapses from 9.2\% to 2.5\%.

\begin{table}[h]
\centering
\caption{Within-window inflation (GraphSAGE, PR-AUC, 5 seeds). The
random-over-temporal ratio cancels each window's base rate.}
\label{tab:windowed}
\begin{tabular}{lcccc}
\toprule
Window & Temporal & Random & Inflation & Ratio \\
\midrule
Pre-shutdown (35--42) & $0.651\pm0.042$ & $0.895\pm0.018$ & $+0.244\pm0.059$ & $1.4\times$ \\
Post-shutdown ($\geq$43) & $0.042\pm0.001$ & $0.348\pm0.033$ & $+0.306\pm0.034$ & $8.2\times$ \\
\bottomrule
\end{tabular}
\end{table}

The absolute difference is directional but not conclusive (paired
$+0.062\pm0.058$, sign-flip $p=0.0625$, 4/5 seeds). On the relative scale the
separation is clear: the within-window random-over-temporal ratio rises from
$1.4\times$ pre-shutdown to $8.2\times$ post-shutdown, and the paired difference
of per-seed log-ratios is $+1.78\pm0.10$ with all five seeds positive (sign-flip
$p=1/32$). Post-shutdown, the temporal model is close to the no-skill floor
(0.042 vs a 0.025 base rate) while the random-split model retains substantial
skill (0.348). This skill can only have been acquired by training on the
post-shutdown period itself. Two caveats: the post-shutdown ratio's denominator
is near the floor, so the ratio's magnitude (though not its sign) is sensitive to
small absolute changes; and the two windows differ in base rate (9.2\% vs 2.5\%),
so cross-window comparisons of \emph{absolute} inflation are imperfect. The
result is consistent with the mechanism's prediction in the relative sense; the
weaker absolute test is reported alongside.

\subsection{Cross-domain controls: a dose-response pattern}
\label{sec:dose}
Two additional datasets bracket the drift axis. On DGraph-Fin
\citep{huang2022dgraph} (3.7M nodes, timestamped edges, temporally stable
$\sim$1.3\% fraud rate), the same GraphSAGE shows essentially no PR-AUC inflation
($+0.001\pm0.001$ over three seeds): where the distribution does not shift, there
is nothing to access. A wider model (hidden dimension 128) behaves identically
(PR-AUC inflation $+0.004$; ROC-AUC $+0.055$). Caveats stated plainly: ROC-AUC
does inflate ($+0.025\pm0.006$; up to $+0.055$ at width 128), both arms sit near
the PR-AUC floor at this model strength, and GADBench reports strong GNNs on
DGraph-Fin, so this control is supporting evidence rather than proof.

On IBM AMLworld HI-Small (5.08M synthetic transactions with complete ground
truth, 0.102\% laundering; tabular XGBoost, three seeds), whose generator drifts
moderately (laundering prevalence 0.075\% $\rightarrow$ 0.152\% across the
stream), the inflation is intermediate: temporal PR-AUC $0.082\pm0.002$ versus
random $0.206\pm0.006$ ($+0.124\pm0.007$). Notably the base-rate effect runs
\emph{against} this inflation (the temporal test window has the \emph{higher}
prevalence), so the $+0.124$ is distribution access net of an adverse base-rate
term. Across four datasets the inflation thus tracks the drift magnitude: stable
(DGraph) $+0.001$; moderate drift (AMLworld) $+0.124$; regime change (Elliptic
family) $+0.41$ to $+0.55$. Three points across different domains and model
classes are a consistent pattern rather than a fitted law, but the dose-response
direction is exactly what the distribution-access mechanism predicts.

\subsection{The unsolved post-shift regime and negative method results}
Under the temporal protocol, a rolling per-window backtest of EvolveGCN-O
\citep{pareja2020evolvegcn} collapses at the dark-market shutdown (steps 44--46:
PR-AUC $\sim$0.01; Figure~\ref{fig:backtest}) and only partially recovers. A
learning-rate $\times$ gradient-clipping sweep (six configurations, seed 42)
improves the aggregate test PR-AUC from 0.100 to at best 0.200 (lr
$5{\times}10^{-3}$, no clipping), still far below every temporal-protocol
baseline; the conclusion is unchanged.

\begin{figure}[t]
\centering
\includegraphics[width=0.75\linewidth]{figs/backtest.pdf}
\caption{Rolling per-window test PR-AUC of EvolveGCN-O under the temporal
protocol. Aggregate scores hide the collapse at the dark-market shutdown
(steps 44--46).}
\label{fig:backtest}
\end{figure}

Is the post-shift regime merely unexplored, or genuinely hard? We tested four
drift-robust variants; none improves on the baselines. A USAD-style two-decoder
adversarial graph autoencoder \citep{audibert2020usad} (naive, rolling-normal,
and domain-adversarial variants) scores PR-AUC $0.0367\pm0.0001$
over three seeds, a stable and decisive negative. A supervised domain-adversarial
GraphSAGE \citep{ganin2015dann} scores $0.462\pm0.090$ over three seeds,
statistically indistinguishable from the plain baseline ($0.510\pm0.039$): the
time-invariance constraint yields no robustness gain. We note a second
self-correction here: at a single seed (42, PR-AUC 0.364) DANN appeared clearly
harmful; the multi-seed run shows that stronger conclusion was seed noise.
Enforcing time-invariance does not solve this regime change; whether it actively
hurts is within noise.

\section{Discussion and recommendations}
(1) On temporally shifting fraud graphs, report illicit-class metrics under a
strict temporal split; random-split numbers can overstate PR-AUC by up to
$+0.55$. (2) The post-shift window is the real benchmark; aggregate scores hide
the collapse. (3) Graph structure is not a free win on Elliptic: features already
encode neighborhood aggregates, and under honest evaluation trees lead. (4)
Method claims should be accompanied by decomposition-style controls; our own
double-leak hypothesis did not survive its ablation.

\section{Limitations}
\label{sec:limitations}
(1) The tuning sweeps and the EvolveGCN sweep are single runs at seed 42; the
headline inflation, decomposition, DGraph control, and negative-method results are
multi-seed (3-5 seeds). (2) The DGraph-Fin control operates near the
PR-AUC floor for the models we trained (PR-AUC 0.036--0.040 against a 1.3\% base
rate), ROC-AUC does inflate there ($+0.025\pm0.006$ over seeds; up to $+0.055$ at
width 128), and GADBench reports far stronger GNNs on this dataset; the control
is therefore supporting evidence for the shift-driven mechanism, not proof. (3)
Our temporal-protocol GNN F1 sits below some honest-protocol literature values
($\sim$0.69 for GraphSAGE) even after tuning; the trees-over-GNNs magnitude we
measure is therefore larger than the field's, though the ordering is consistent.
(4) The trees-over-GNNs ordering is established here on Elliptic-family data
only. (5) The GAT equivalence bound on message-passing leakage is loose
($\pm 0.05$). (6) DGraph-Fin's temporal stability is measured coarsely (labeled
fraud prevalence 1.2\% in the train window versus 1.4\% in the test window under
our first-seen split), not characterized in depth. (7) Per concurrent work
\citep{safar2026opacity}, Elliptic's features embed construction-level look-ahead
(address aggregates computed after the test period) and entity nonindependence;
our ``honest'' protocol is honest at the \emph{split} level only. Both arms share
the same features, so exposure is symmetric; the induced effect on each arm need
not be, so construction-level leakage primarily qualifies our absolute
temporal-arm scores (upper bounds on deployable performance) and could
second-order perturb the arm differences on which the decomposition rests.

\section{Conclusion}
The near-perfect reported performance on Elliptic-family benchmarks is largely an
evaluation artifact whose dominant mechanism is distribution access, not
graph-specific leakage. We provide the decomposition, the controls, a reproducible
harness together with negative results reported in full, and we re-anchor the
state of the art under the temporal protocol.

\section*{Reproducibility statement}
Every number in this paper is produced by one seeded CLI command. The anonymized
supplementary archive contains the full source code, the per-seed values for
every cell, and the statistical summary
(\texttt{docs/results/stats\_summary.json}). The public repository, which
additionally carries a continuous-integration pipeline (lint, typecheck, tests,
Docker smoke) on every commit, is linked in the camera-ready version. Elliptic
downloads automatically via PyTorch Geometric; Elliptic++ and DGraph-Fin are
obtained from their authors (DGraph-Fin is gated and requires registration);
AMLworld is obtained from Kaggle.

\bibliographystyle{unsrtnat}
\bibliography{references}

\appendix

\section{Per-seed values for the headline inflation}
\label{app:perseed}
Table~\ref{tab:perseed} lists the per-seed PR-AUC values behind
Table~\ref{tab:inflation}; the released results files contain the same
per-seed detail for every other cell in the paper.

\begin{table}[h]
\centering
\caption{Per-seed illicit-class PR-AUC (seeds 42--46).}
\label{tab:perseed}
\small
\begin{tabular}{llccccc}
\toprule
Model & Arm & s42 & s43 & s44 & s45 & s46 \\
\midrule
\multirow{2}{*}{GCN} & temporal & 0.294 & 0.324 & 0.206 & 0.201 & 0.302 \\
 & random & 0.800 & 0.820 & 0.824 & 0.822 & 0.823 \\
\multirow{2}{*}{GraphSAGE} & temporal & 0.488 & 0.542 & 0.455 & 0.546 & 0.522 \\
 & random & 0.925 & 0.912 & 0.930 & 0.913 & 0.922 \\
\multirow{2}{*}{GAT} & temporal & 0.332 & 0.251 & 0.285 & 0.236 & 0.384 \\
 & random & 0.811 & 0.783 & 0.746 & 0.730 & 0.806 \\
\multirow{2}{*}{GraphTransformer} & temporal & 0.487 & 0.436 & 0.498 & 0.475 & 0.540 \\
 & random & 0.897 & 0.899 & 0.912 & 0.909 & 0.896 \\
\multirow{2}{*}{XGBoost} & temporal & 0.790 & 0.790 & 0.793 & 0.792 & 0.790 \\
 & random & 0.987 & 0.984 & 0.986 & 0.986 & 0.984 \\
\bottomrule
\end{tabular}
\end{table}

\section{GenAI usage statement}
Code, experiments, and manuscript drafting were carried out with substantial
assistance from an LLM-based coding agent (Claude, Anthropic), with the agent
workflow documented in the repository. The research questions, methodological
decisions, and final claims were set and verified by the human author, who takes
full responsibility for the content.

\end{document}
